\section{Limitations and Threats to Validity}

\textbf{Empirical scope.}  The study deliberately covers one H100, one 7B
model, one vLLM version, BF16, fixed output length, and bounded
workload/configuration grids.  The measured knob search varies sequence and
token limits plus chunked prefill.  The claims are therefore about evidence
allocation on one target GPU, not parallel or distributed serving.

\textbf{Fidelity and telemetry.}  DCGM cumulative energy captures GPU-device
rather than host or facility energy.  The CPU Sandbox and short Probe may both
miss behavior outside the frozen model, software, workload, and hardware scope;
the Sandbox TTFT inversion observed here further shows that good energy ranking
must not be generalized to every metric.  TP/PP, cross-hardware, MoE, and
multi-node communication require new measured calibration and verification and
are reserved for a broader follow-on study.

\textbf{Uncertainty and policy value.}  Prediction intervals are not calibrated
across hardware generations; new kernels, engines, GPUs, arrival processes, or
MoE routing may invalidate residual corrections.  The current \ipig{} estimator
is an auditable proxy rather than exact Pareto mutual information.  Its
\PolicyEpisodes{} episodes per policy bootstrap only \PolicyHardwareRepeats{}
real repeats per candidate--stage key, so they measure sensitivity to observed
run-to-run variation rather than independent production deployments.  \ipig{}
beats random and cost-blind acquisition on the declared cost--success metrics
but not cheapest-first in this small grid.  The exhaustive
\ConfigurationTotalGPUHours{}-GPU-h corpus validates fidelity and supplies an
oracle; replay does not erase that offline data-collection cost.
